\section{Homogeneous chains in momentum space}
\label{sec:fourier}
\motiv{When every site of the periodic chain of Section~\ref{sec:chain} carries the same
on-site covariance $\Gamma$, the chain covariance commutes with the shift of sites, and
the discrete Fourier transform over sites splits it into $N$ blocks of size $2P\times2P$,
one per momentum $k$. Each block is again a contraction (Definition~\ref{def:contract}),
now of a single site against a single bond, with the shift of the right virtual modes
turned into a phase, and the inverse transform recovers $\Gamma_N$ from its blocks;
antisymmetry and admissibility of $\Gamma$ pass to every block.}

Throughout this section $N\geq1$ and $M,P\geq0$; the on-site covariances are
$\Gamma_i=\Gamma$ for every site, a fixed real matrix on $I_{\mathrm{site}}=(x_A,x_B,x_C,y_A,y_B,y_C)$, and
$\Gamma_N=\Gamma_N(\Gamma,\ldots,\Gamma)$ is the covariance of the chain
(Definition~\ref{def:periodic}), with $D$ the $E\times E$ block of $\Gamma_{\mathrm{in}}$.
We write retained labels as $(i,a)$ with $i\in\mathbb Z/N$ the site and
$a\in R_0=(x_C,y_C)$ one of the $2P$ physical labels of a site. The site labels split as
$I_{\mathrm{site}}=R_0\sqcup E_0$ with $E_0=(x_A,x_B,y_A,y_B)$; $\beta_M$ is the real bond covariance of
$E_0$ (Definition~\ref{def:bond} with $L=A$, $L'=B$ and the index set $\{1,\ldots,M\}$).
For a complex matrix $X$, $X^*$ is the conjugate transpose, $\overline X$ the entrywise
conjugate, and $X^{-1}$ the inverse of an invertible $X$.

\subsection{Twisted on-site covariances}
\begin{definition}[Characters and the twisted on-site covariance]\label{def:twist}
For $k,i\in\mathbb Z/N$ let $\chi_k(i)=e^{-2\pi\ii\,ki/N}$ and $\omega_k=\chi_k(1)$.
The \emph{twist} $t_k:I_{\mathrm{site}}\to\C$ is $t_k(u)=\omega_k$ for $u\in\{x_B,y_B\}$ and
$t_k(u)=1$ otherwise, and the \emph{twisted on-site covariance} at momentum $k$ (the on-site covariance conjugated by
the momentum phase on the right virtual modes) is
\[
 \Gamma^{(k)}=\operatorname{diag}(t_k)\,\Gamma\,\operatorname{diag}(t_k)^*,\qquad
 \Gamma^{(k)}_{uv}=t_k(u)\,\Gamma_{uv}\,\overline{t_k(v)} .
\]
\end{definition}
With this sign of $\chi_k$, Theorem~\ref{thm:reconstruct} expands $\Gamma_N$ in
$e^{+2\pi\ii\,k(i-j)/N}\,\Sigma(k)$. \vet{The opposite convention expands in} $e^{-2\pi\ii\,k(i-j)/N}$, and
in that convention the block at momentum $k$ is $\Sigma(-k)=\overline{\Sigma(k)}$
(Theorem~\ref{thm:blocks}(a)).

\begin{definition}[Momentum-space covariance]\label{def:symbol}
Write $\Gamma^{(k)}=\begin{pmatrix}A_k&B_k\\C_k&D_k\end{pmatrix}$ in the split
$I_{\mathrm{site}}=R_0\sqcup E_0$. The \emph{momentum-space covariance} at $k$ (the symbol of the
block-circulant matrix $\Gamma_N$) is the complex $R_0\times R_0$ matrix
\[
 \Sigma(k)=A_k+B_k\,\beta_M\,(1-D_k\beta_M)^{-1}\,C_k ,
\]
the contraction formula of Definition~\ref{def:contract} applied over $\C$ to $\Gamma^{(k)}$
and the bond of one site.
\end{definition}
\artefact{When $1-D_k\beta_M$ is singular, we read $\Sigma(k)$ with $(1-D_k\beta_M)^{-1}=0$,
the value of the formalization's inverse on singular matrices; only
Theorem~\ref{thm:blocks}(a) uses $\Sigma(k)$ in that case.}

\begin{definition}[Fourier coefficients of the chain]\label{def:fourier}
For $k,l\in\mathbb Z/N$ and $a,b\in R_0$,
\[
 \widehat{\Gamma}_N(k,l)_{ab}=\frac1N\sum_{i,j\in\mathbb Z/N}\chi_k(i)\,(\Gamma_N)_{(i,a)(j,b)}\,\overline{\chi_l(j)} .
\]
\end{definition}

\subsection{The Fourier transform of the chain}
Let $\Phi_N$ be the $N\times N$ matrix with entries $(\Phi_N)_{ki}=\chi_k(i)/\sqrt N$; by the
orthogonality of characters, $\sum_i\chi_k(i)\overline{\chi_l(i)}=N\delta_{kl}$, it is
unitary. Identify the retained labels with pairs, $R\cong(\mathbb Z/N)\times R_0$ by
$(i,a)$, and the eliminated labels likewise, $E\cong(\mathbb Z/N)\times E_0$ by
$x^{L}_{(i,a)}\mapsto(i,x_{A,a})$, $x^{L'}_{(i,a)}\mapsto(i,x_{B,a})$ and the same for $y$;
a label $(i,u)$ with $u\in\{x_B,y_B\}$ is thus the right virtual mode of site $i-1$
(Definition~\ref{def:input}). Let $\Phi_R=\Phi_N\otimes1_{R_0}$,
$\Phi_E=\Phi_N\otimes1_{E_0}$ and $\Phi=\Phi_R\oplus\Phi_E$ on $R\sqcup E$.

\begin{lemma}[Block diagonalization of the input]\label{lem:input-blocks}
$\Phi_E\beta_N\Phi_E^*=\beta_N$, and $\Phi\Gamma_{\mathrm{in}}\Phi^*=\bigoplus_{k}\Gamma^{(k)}$ in the
split $R\sqcup E$, that is, the four blocks of $\Phi\Gamma_{\mathrm{in}}\Phi^*$ are
$\bigoplus_kA_k$, $\bigoplus_kB_k$, $\bigoplus_kC_k$, $\bigoplus_kD_k$.
\end{lemma}
\begin{proof}
\begin{enumerate}
\item[(a)] In the pair identification, $\beta_N=1_N\otimes\beta_M$, because the bond
pairs $x^{L}_{(i,a)}$ with $x^{L'}_{(i,a)}$ at the same pair $(i,a)$
(Definition~\ref{def:periodic}); it commutes with $\Phi_N\otimes1$.
\item[(b)] Let $s(u)=1$ for $u\in\{x_B,y_B\}$ and $s(u)=0$ otherwise. The label $(i,u)$
with $u\in\{x_B,y_B\}$ names a mode of site $i-1$, and $\operatorname{diag}(\Gamma_\bullet)$
is site diagonal, so
$(\Gamma_{\mathrm{in}})_{(i,u)(j,v)}=\Gamma_{uv}\,[\,i-s(u)=j-s(v)\,]$. Hence
\[
 (\Phi\Gamma_{\mathrm{in}}\Phi^*)_{(k,u)(l,v)}
 =\frac{\Gamma_{uv}}N\sum_{i}\chi_k(i)\,\overline{\chi_l(i-s(u)+s(v))}
 =\delta_{kl}\,\Gamma_{uv}\,\omega_k^{\,s(u)}\,\overline{\omega_k}^{\,s(v)}
 =\delta_{kl}\,\Gamma^{(k)}_{uv},
\]
by $\chi_l(i+m)=\chi_l(i)\chi_l(m)$, $\chi_l(1)=\omega_l$, and orthogonality; the last
equality is Definition~\ref{def:twist}. \qedhere
\end{enumerate}
\end{proof}

\begin{lemma}[Contraction and relabelling]\label{lem:contract-conj}
Let $\Gamma$ be a matrix on $R\sqcup E$, $\Gamma_p$ a matrix on $E$ with $1-D\Gamma_p$
invertible, and $Y_R$, $Y_E$ invertible matrices on the retained and on the eliminated
labels. Then, with the contraction formula taken over $\C$,
\[
 \begin{pmatrix}Y_RAY_R^{-1}&Y_RBY_E^{-1}\\Y_ECY_R^{-1}&Y_EDY_E^{-1}\end{pmatrix}\star(Y_E\Gamma_pY_E^{-1})
 =Y_R\,(\Gamma\star\Gamma_p)\,Y_R^{-1} .
\]
If $\Gamma=\bigoplus_k\Gamma_k$ and $\Gamma_p=\bigoplus_k\Gamma_{p,k}$ are block diagonal
in a common index $k$, then $\det(1-D\Gamma_p)=\prod_k\det(1-D_k\Gamma_{p,k})$, and when
this is nonzero, $\Gamma\star\Gamma_p=\bigoplus_k\Gamma_k\star\Gamma_{p,k}$.
\end{lemma}
\begin{proof}
$1-Y_EDY_E^{-1}\,Y_E\Gamma_pY_E^{-1}=Y_E(1-D\Gamma_p)Y_E^{-1}$, whose inverse is
$Y_E(1-D\Gamma_p)^{-1}Y_E^{-1}$; the four products then telescope. For block
diagonal data, determinants multiply over blocks, the inverse of an invertible block
diagonal matrix is block diagonal with the inverse blocks, and sums and products are
blockwise.
\end{proof}

\begin{theorem}[Determinant factorization]\label{thm:det-factor}
$\det(1-D\beta_N)=\prod_{k\in\mathbb Z/N}\det(1-D_k\beta_M)$. In particular, if
$1-D\beta_N$ is invertible then $1-D_k\beta_M$ is invertible for every $k$.
\end{theorem}
\begin{proof}
$\Phi_E(1-D\beta_N)\Phi_E^*=1-\bigoplus_kD_k\beta_M$ by Lemma~\ref{lem:input-blocks};
$\Phi_E$ is unitary, so the determinants agree, and the block determinant is the product
(Lemma~\ref{lem:contract-conj}). A nonzero product has nonzero factors.
\end{proof}

\begin{theorem}[Translation invariance and diagonalization]\label{thm:diagonal}
Let $1-D\beta_N$ be invertible. Then for all $i,j\in\mathbb Z/N$, $k,l\in\mathbb Z/N$
and $a,b\in R_0$,
\[
 (\Gamma_N)_{(i+1,a)(j+1,b)}=(\Gamma_N)_{(i,a)(j,b)},
 \qquad
 \widehat{\Gamma}_N(k,l)=\begin{cases}\Sigma(k)&k=l,\\0&k\neq l.\end{cases}
\]
\end{theorem}
\begin{theorem}[Reconstruction]\label{thm:reconstruct}
Let $1-D\beta_N$ be invertible. Then for all $i,j\in\mathbb Z/N$ and $a,b\in R_0$,
\[
 (\Gamma_N)_{(i,a)(j,b)}=\frac1N\sum_{k\in\mathbb Z/N}\overline{\chi_k(i)}\,\Sigma(k)_{ab}\,\chi_k(j).
\]
\end{theorem}
\begin{proof}[Proof of Theorems~\ref{thm:diagonal} and~\ref{thm:reconstruct}]
\begin{enumerate}
\item[(a)] By Lemma~\ref{lem:input-blocks}, Lemma~\ref{lem:contract-conj} with
$Y=\Phi$ (so $Y^{-1}=\Phi^*$), and Theorem~\ref{thm:det-factor},
\[
 \Phi_R\,\Gamma_N\,\Phi_R^*=(\Phi\Gamma_{\mathrm{in}}\Phi^*)\star(\Phi_E\beta_N\Phi_E^*)
 =\Bigl(\bigoplus_k\Gamma^{(k)}\Bigr)\star\Bigl(\bigoplus_k\beta_M\Bigr)=\bigoplus_k\Sigma(k).
\]
\item[(b)] The $((k,a),(l,b))$ entry of the left side is $\widehat{\Gamma}_N(k,l)_{ab}$
by the definition of $\Phi_R$; this is the second part of Theorem~\ref{thm:diagonal}.
\item[(c)] Multiplying (a) by $\Phi_R^*$ on the left and $\Phi_R$ on the right gives
$\Gamma_N=\Phi_R^*\bigl(\bigoplus_k\Sigma(k)\bigr)\Phi_R$, whose $((i,a),(j,b))$ entry is the
sum of Theorem~\ref{thm:reconstruct}.
\item[(d)] $\overline{\chi_k(i+1)}\,\chi_k(j+1)=\overline{\chi_k(i)}\,\overline{\omega_k}\,\chi_k(j)\,\omega_k=\overline{\chi_k(i)}\,\chi_k(j)$,
so replacing $(i,j)$ by $(i+1,j+1)$ leaves the sum in (c) unchanged; this is the first
part of Theorem~\ref{thm:diagonal}. \qedhere
\end{enumerate}
\end{proof}
\begin{remark}\label{rem:singular}
The invertibility hypothesis enters translation invariance only through this route; the
chain construction is shift symmetric, and we make no claim here about the singular case.
\end{remark}

\subsection{Admissibility in momentum space}
\begin{theorem}[Symmetries and admissibility of the momentum-space covariance]\label{thm:blocks}
\begin{enumerate}
\item[(a)] If $1-D_k\beta_M$ is invertible then $\Sigma(-k)=\overline{\Sigma(k)}$;
\artefact{the \revs{formalized} statement holds for every $k$, with the inverse of a singular matrix read as $0$.}
\item[(b)] If $\Gamma^{\T}=-\Gamma$ and $1-D\beta_N$ is invertible, then
$\Sigma(k)^*=-\Sigma(k)$ for every $k$.
\item[(c)] If $\Gamma$ is admissible (Definition~\ref{def:admissible}) and $1-D\beta_N$
is invertible, then $1-\Sigma(k)^*\Sigma(k)\succeq0$ for every $k$, that is,
$\|\Sigma(k)z\|\leq\|z\|$ for every $z\in\C^{R_0}$.
\end{enumerate}
\end{theorem}
\begin{proof}
\begin{enumerate}
\item[(a)] $\omega_{-k}=\overline{\omega_k}$, so $t_{-k}=\overline{t_k}$ and, $\Gamma$
being real, $\revs{\Gamma^{(-k)}}=\overline{\Gamma^{(k)}}$; $\beta_M$ is real, and the inverse, read as
$0$ on singular matrices, commutes with entrywise conjugation, so
the formula of Definition~\ref{def:symbol} commutes with conjugation.
\item[(b)] $\Gamma_N^{\T}=-\Gamma_N$ (Theorem~\ref{thm:antisym}, applied to
$\Gamma_{\mathrm{in}}$, which is antisymmetric with $\Gamma$, and to $\beta_N$), and
$\Gamma_N$ is real, so $\Gamma_N^*=-\Gamma_N$ and $\Phi_R\Gamma_N\Phi_R^*$ is skew-Hermitian;
by step (a) of the proof of Theorem~\ref{thm:diagonal} its blocks $\Sigma(k)$ are skew-Hermitian.
\item[(c)] $\Gamma_N$ is admissible (Theorem~\ref{thm:periodic}), so
$1-\Gamma_N^{\T}\Gamma_N\succeq0$ over $\R$, hence over $\C$: for $z=a+\ii b$ with $a,b$
real, $z^*Xz=a^{\T}Xa+b^{\T}Xb$ when $X$ is real symmetric. Conjugating by $\Phi_R$ and
using step (a) of the proof of Theorem~\ref{thm:diagonal},
$\Phi_R(1-\Gamma_N^{\T}\Gamma_N)\Phi_R^*=1-\bigoplus_k\Sigma(k)^*\Sigma(k)\succeq0$, and a
block diagonal matrix is positive semidefinite exactly when each block is (test with
vectors supported on one block). \qedhere
\end{enumerate}
\end{proof}
\motiv{Part (c) is the momentum-space form of admissibility. For pure $\Gamma$ the chain
covariance $\Gamma_N$ is pure (Corollary~\ref{cor:chain-pure}, a matrix fact whose proof in these notes
goes through Fock space), and the argument of (c) with
$1-\Gamma_N^{\T}\Gamma_N=0$ gives $\Sigma(k)^*\Sigma(k)=1$ for every $k$.}
